\section{Background}

\paragraph{Roofline.} The roofline model bounds attainable performance by
$\min(\text{peak compute}, \text{peak bandwidth} \times \text{arithmetic
intensity})$~\cite{williams2009}. The ridge, peak compute over peak bandwidth,
separates the memory bound region (left) from the compute bound region (right).
GEMM at large $N$ has an arithmetic intensity that grows with $N$, so it lives far
to the right of the ridge; a slow GEMM is therefore an implementation gap, not an
algorithmic ceiling.

\paragraph{Blocking theory.} The decisive step from a shared memory tiled kernel
to a fast one is register blocking: each thread computes many outputs, which lifts
the arithmetic intensity per thread and the work done per shared memory
instruction. This follows Volkov and Demmel~\cite{volkov2008}, and the diagnostic
rounds below confirm it on Blackwell.

\paragraph{Tensor pipeline.} The fifth generation tensor cores raise the compute
ceiling far above the FP32 units. Reaching that ceiling means feeding the cores
fast enough, which is a shared memory problem: the fragment loads, not the matrix
multiply instructions, are what saturate first.
